\subsection{Patrones de diseño y continuidad}

\label{sec:continuidad-logica}

La continuidad se articula con ISO 22301 para la gestión de continuidad del negocio y con ISO/IEC 27031 para la preparación de las TIC que la sostiene, conforme a RT-10.03 y RT-10.04 (ISO, 2019, 2025). En esta vista, su aplicación se concreta en las dependencias entre módulos, las funciones que pueden continuar desconectadas, los estados sin confirmar y las condiciones para volver a una operación consistente. Esta correspondencia de diseño no constituye una certificación ni acredita por sí sola el cumplimiento integral de las normas.

Ante pérdida de conectividad, M1, M2 y M5 conservan la operación local y M9 mantiene el bloqueo sanitario; M3 y M6 conservan la captura de terreno durante las autonomías definidas. Las confirmaciones que dependen de crédito, stock central o documentos del ERP permanecen explícitamente sin confirmar o bloqueadas según su regla. La continuidad protege la evidencia y la seguridad de la operación: no habilita una salida sin autorización tributaria ni permite levantar un bloqueo sanitario para recuperar velocidad.

La recuperación lógica exige restablecer identidad y autorización, disponer del estado transaccional y de los eventos sin confirmar, y habilitar después el procesamiento y la reconciliación. Cada consumidor confirma solo operaciones persistidas, deduplica los reenvíos y deriva los conflictos a su responsable. Antes de cerrar la contingencia se cotejan pedidos, reservas, lotes, entregas y rendiciones con sus identificadores de origen; las operaciones sin confirmar y las excepciones quedan visibles y auditadas. El orden concreto de restauración y la capacidad de respaldo se documentan en 4.2 y 4.3.

La evidencia de aceptación vincula cada servicio con su requisito, módulo, dependencia, objetivo de recuperación y prueba. Se ensayan pérdida de WAN durante 24 horas en el sitio, captura móvil durante 14 horas, indisponibilidad del ERP, relevo sin IdP y recuperación con mensajes repetidos. Se mide la conservación de operaciones confirmadas, el respeto de los bloqueos, la ausencia de duplicados y el tiempo hasta recuperar un estado conciliado. Una pérdida completa del sitio requiere una prueba de restauración distinta del corte WAN, contrastada con RTO $\leq$ 4 horas y RPO $\leq$ 15 minutos para los servicios críticos. La extracción continua por fibra, LTE y satélite sostiene el RPO, verificado con AL-DR-01; 4.3.2 justifica el límite residual.

La arquitectura aplica un conjunto de patrones de diseño que responden a las restricciones específicas de la operación de Puelche:

\textbf{Registro local y sincronización idempotente.} Cada escritura recibe un UUID declarado por el cliente que se conserva al reintentar. El evento permanece en la cola local hasta que el servicio de negocio confirma su procesamiento. Ese servicio deduplica dentro de una ventana documentada (RT-02.06) y resuelve conflictos con reglas deterministas y bitácora auditable (RT-03.12). La copia de lectura en SQLite/Room se actualiza por separado. Así, renovar precios o stock no elimina un pedido sin confirmar y el trabajo de terreno no depende de una consulta a Redis en tiempo real.

\textbf{Servicios sin estado durable en memoria.} Las bases de datos y las colas conservan el estado de negocio; las sesiones utilizan los mecanismos de identidad definidos. Sustituir una instancia no debe perder el progreso de una operación. El acceso desconectado respeta las vigencias de las credenciales y las restricciones explicadas en la capa de seguridad.

\textbf{Degradación controlada sin pérdida silenciosa.} Cuando un componente no responde, la operación continúa en modo reducido informado a la persona usuaria ("modo offline"). Ninguna degradación produce pérdida de una transacción de venta o entrega: si una escritura no llega al servidor, queda en el buffer local del dispositivo, visible como no confirmada y reconciliada en la sincronización posterior (RT-02.09). La clasificación de servicios (RT-10.02) distingue cuatro niveles, conforme a la Tabla~\ref{tab:jd05}: crítico, preparación, despacho y emisión de la guía durante la ventana 05:30--07:00, con el bloqueo por excursión térmica y la identidad de bodega que lo habilitan; alto, toma de pedido y consulta de stock y crédito, entrega y evidencia, planificación de rutas, los demás DTE, integración con el ERP, cobranza y rendición, registro de temperatura; medio, canal moderno (EDI), notificaciones, analítica y tableros, portales, consultas de geolocalización, registros de auditoría y seguridad; bajo, reportería e informes a pedido.

\textbf{Resiliencia con cortacircuitos y colas.} Las llamadas remotas tienen un tiempo máximo de espera explícito (RT-02.08). Las operaciones reintentables utilizan espera creciente con variación aleatoria y un cortacircuitos que suspende temporalmente las llamadas a un servicio que falla. Si el ERP no está disponible, la preventa conserva la captura de pedidos, mientras las confirmaciones dependientes del legado permanecen sin confirmar. Las colas y sus DLQ permiten recuperar el procesamiento; los consumidores deben deduplicar porque un mensaje puede entregarse más de una vez.

\textbf{Integración con legados mediante ACL.} Los módulos M1, M5, M7, M8 y M11 intercambian información con el ERP a través de adaptadores; no absorben su responsabilidad tributaria ni modifican su código. Las capacidades del WMS 2013 se reemplazan gradualmente en M1, M2 y M5. La ACL mantiene separados el modelo de negocio de la solución y los formatos del legado, en coherencia con RT-02.14. No se permiten escrituras directas al ERP.

\textbf{Correlación de extremo a extremo.} Cada transacción lleva un \texttt{transaction\_id} que acompaña las solicitudes y los eventos, incluidos los capturados sin conexión. Al sincronizar, ese identificador permite relacionar métricas, trazas y registros en la plataforma común de observabilidad de nube y sitios locales (RT-03.16 y Art. 16.4).

\textbf{Gobierno de contratos y versionado.} Las APIs de negocio y de integración se documentan desde el código, con actualización automática de OpenAPI 3.1 (síncrona) y AsyncAPI 2.6 (eventos), semver estricto, propietario declarado por módulo, compatibilidad hacia atrás y aviso mínimo de 6 meses antes de deprecar una versión (RT-05.16/17).

\textbf{Seguridad Zero Trust.} Cada solicitud se verifica explícitamente, sin presuponer la confiabilidad de ninguna red, dispositivo o identidad (NIST, 2020, SP~800-207). Los servicios utilizan mTLS (RT-05.18), los secretos tienen rotación automática (RT-04.09) y los registros de auditoría son inmutables (RT-16.07). La segmentación y las conexiones entrantes al sitio local se especifican en 4.2.

\textbf{Separación transaccional y analítica.} Los tableros consultan Redshift Serverless, no el transaccional en vivo. Los eventos incrementales permiten mantener las latencias de hasta 5 minutos para operación, 2 horas para cierre comercial y 4 horas para gestión. Las tareas pesadas se programan fuera de la ventana crítica, sin detener la actualización incremental necesaria para el tablero operacional.

\subsubsection{SOLID aplicado a los contratos de negocio}

La responsabilidad única se comprueba en el flujo de un pedido: M3 captura y confirma el pedido, M2 decide si reserva stock, M5 ejecuta la preparación y la ACL traduce hacia el ERP. Ninguno de esos componentes puede asumir la decisión del otro por compartir un proceso Laravel. Para extender el canal moderno, M11 incorpora un adaptador por cadena bajo un contrato publicado; M3 no se modifica para aprender un nuevo formato EDI. Esa es la aplicación concreta del principio abierto/cerrado.

Un adaptador nuevo sustituye a uno previo solo si conserva las respuestas, errores e idempotencia exigidos por el contrato versionado: las pruebas de contrato verifican esa sustitución. Las interfaces se separan por capacidad y autorización ---consulta de stock, reserva, recepción de POD y liberación sanitaria--- para que un cliente de lectura no reciba por accidente una operación de escritura. Finalmente, M3, M5 y M7 dependen de puertos de negocio para inventario, documentos y mensajería, no de llamadas directas al SDK de una nube o a tablas del ERP. En QA se reemplazan esos puertos por dobles de prueba; en producción los adaptadores implementan el contrato. Las pruebas de contrato verifican sustitución, segregación de interfaces e inversión de dependencias.
